\begin{abstract}
\NHFOAM\ is a non-hydrostatic, free-surface wave model for OpenFOAM in the
spirit of NHWAVE and SWASH. It solves the incompressible Navier--Stokes
equations on surface- and terrain-following ($\sigma$) meshes that move with the
free surface and can be unstructured in the horizontal. As in NHWAVE, the
pressure is stored on the layer interfaces, with one surface unknown per
column that carries the free surface implicitly; the resulting matrix is
solved with conjugate gradients and a column (line) preconditioner. The model
has relaxation zones and a boundary wavemaker for wave generation and absorption, prescribed seabed
motion, and depth-induced breaking with the hydrostatic front approximation
and a breaking eddy viscosity.
This report documents the equations and the numerical method, and collects
the verification and validation to date. Run on the same set-ups as NHWAVE,
\NHFOAM\ resolves dispersion to $kh = 7$ with two layers (period error below
0.6\%), and over a submerged bar and for a 3D landslide it matches the
measurements as well as NHWAVE and REEF3D::NHFLOW. The tests of
\citet{ma2012} were then repeated against NHWAVE from the simplest upwards,
and they exposed two weaknesses that are now fixed. The convection term was
linearised with the flux of the previous time level, which made the time
integration first order for nonlinear waves: a solitary wave of $H/h = 0.5$
lost 6.4\% of its height in 150 m. Extrapolating the convecting flux restores
second order at no cost and brings the loss to 1.0\%, against 2.3\% for NHWAVE.
And the TVD convection scheme damped short standing waves; with filtered
central differencing the standing-wave error is 0.07--0.56\%, against
0.6--4.8\% for NHWAVE. On the 3D elliptic shoal of \citet{berkhoff1982},
\NHFOAM\ reproduces the measured focusing on a grid four times coarser than
the paper's, with errors at the first sections close to those of NHWAVE on
the paper's grid, while NHWAVE on the coarse grid loses 30\% of the incident
wave height. Wetting and drying was then added, after NHWAVE's thin-film mask
but adapted to the implicit free surface: a limiter on the outflows keeps the
water depth positive and the volume conserved, and the side fluxes carry the
upwind (TVD) water depth and, in supercritical flow, the upwind predicted
flux. Still water stays at rest against a dry beach, the planar surface in
the parabolic basin of \citet{thacker1981} is reproduced to 0.9 mm rms, and
the breaking solitary wave of \citet{synolakis1987} runs up to $R/d = 0.531$,
against 0.512 for NHWAVE and 0.49--0.55 in the experiments. On the spilling
breakers of \citet{ting1994}, with stream-function wave generation, shoaling
and the breaking point agree with the measurements, but in the surf zone the
hydrostatic front approximation alone, like NHWAVE, loses too little energy.
A persistence threshold that falls with the age of the breaking event and
the breaking eddy viscosity of \citet{kennedy2000} nearly halve the surf-zone
errors on this beach and on that of \citet{hansen1979} (mean error of the
wave height $+6.1 \to +3.5$ cm, NHWAVE $+4.8$ cm) and keep the measured shape
of the undertow. With them, the longshore current of \citet{visser1991}, the
last test of \citet{ma2012}, was run in 3D with periodic lateral boundaries,
with both models on the same coarse grid. Both put the current at about the
measured position and give two thirds of its strength (\NHFOAM\ 0.27 m/s,
NHWAVE 0.25 m/s, measured 0.39 m/s); with the hydrostatic front alone the
current is weaker and too close to the shoreline. The slow decay of the wave
height just after breaking is now the main limitation: no parameter of the
breaking model removes it, although a kinematic onset with the crest speed
estimated from the wave face breaks at the measured place and lowers the
surf-zone error to $+2.4$ cm, and a stabilised $k$--$\omega$ closure gives
the undertow its measured strength but does not replace the breaking model.
A boundary wavemaker with active absorption, a return current for the
Stokes drift and irregular (JONSWAP, directional) waves complete the wave
input. The return current removes the slow drift of long runs over the
submerged bar, after which \NHFOAM's errors there are steady and its phase
agrees with that of NHWAVE and REEF3D::NHFLOW. On triangulated meshes the
breaking solitary wave runs up to $R/d = 0.49$ and the waves diffracted by
a vertical cylinder are within 7\% of linear theory. Around the conical
island of \citet{briggs1995}, the first three-dimensional test of wetting and
drying, \NHFOAM\ and NHWAVE on the same grid agree to within 4\% at most
gauges and to $0.02d$ on average in the run-up, and both reproduce the shape
of the measured run-up around the island; \NHFOAM\ was 3 to 6 times faster.
Measured gauge and paddle records can now drive the boundary wavemaker: a
solitary wave made from its own record arrives within 0.5\% of its height
and 10 ms of the exact wave after 12 m up to $H/h = 0.3$, and on the
composite beach of NTHMP benchmark 5 the record of the incident gauge
reproduces the measured incident crests within 10\% in the two smaller
cases, while the laboratory's paddle record gives waves 7--27\% too high.
On that beach \NHFOAM\ and NHWAVE agree within 2--6\% with the same input
and grid, and NHWAVE's $k$--$\varepsilon$ model damps the steeper waves
before they reach the slope. In the Monai valley of NTHMP benchmark 7, where
a tsunami ran up a narrow valley, the benchmark's incident wave drives the
boundary wavemaker directly: on the same grid the two models flood the same
ground and follow the measured gauges to 0.5--0.6 cm rms, NHWAVE's bore
arriving 0.3--0.4 s late, and on the benchmark grid \NHFOAM's run-up at the
head of the valley is 7.4 cm against 8.0--10.0 cm measured. Restarts are
exact, parallel runs reproduce serial ones with breakers across processor
boundaries, and a regression suite guards both.

\vspace{0.5em}
\noindent\textbf{Keywords:} non-hydrostatic model; $\sigma$-coordinate;
OpenFOAM; NHWAVE; wave dispersion; wetting and drying; wave run-up;
wave breaking; longshore current; wave generation; active absorption;
unstructured mesh; model validation
\end{abstract}
